\clearpage
\folio{53}%
\chapterhead{4}{IV}{Ideal classes and the unit theorem}{IDEAL CLASSES AND THE UNIT THEOREM}
\renewcommand{\chaptertitle}{IDEAL CLASSES AND THE UNIT THEOREM}

The present chapter is devoted to two important finiteness theorems. For the
proofs of these theorems we shall need some analysis. We proceed to develop
this material first.

%% ----------------------------------------------------------------------
\sechead{1}{Preliminaries concerning discrete subgroups of $\R^{n}$}%
        {PRELIMINARIES CONCERNING DISCRETE SUBGROUPS OF $\R^{n}$}

A subgroup $H$ of $\R^{n}$ is discrete if and only if, for any compact
subset $K$ of $\R^{n}$, the intersection $H\cap K$ is finite. A typical
example of a discrete subgroup of $\R^{n}$ is $\Z^{n}$. We are going to show
that it is almost the only one.

\begin{theorem}
\label{thm:IV.1.1}
Let $H$ be a discrete subgroup of $\R^{n}$. Then $H$ is generated (as a
$\Z$-module) by $r$ vectors which are linearly independent over $\R$ (so
$r\leq n$).
\end{theorem}

\begin{proof}
Let $(e_{1},\ldots,e_{r})$ be a set of elements of $H$ which are linearly
independent over $\R$, where $r$ is as large as possible (again, $r\leq n$).
Let
\begin{equation}
P=\Bigl\{x\in\R^{n}\mid x=\sum_{i=1}^{r}\alpha_{i}e_{i},\quad
0\leq\alpha_{i}\leq1\Bigr\},
\label{eq:IV.1.1}
\end{equation}
the parallelotope constructed on these vectors. Clearly, $P$ is compact, so
$P\cap H$ is finite. Take $x\in H$. From the maximality of the set
$(e_{1},\ldots,e_{r})$ it follows that $x=\sum_{i=1}^{r}\lambda_{i}e_{i}$,
$\lambda_{i}\in\R$. For $j\in\Z$ set
\begin{equation}
x_{j}=jx-\sum_{i=1}^{r}[j\lambda_{i}]e_{i}
\label{eq:IV.1.2}
\end{equation}
(where $[\mu]$ denotes the largest integer less than or equal to $\mu\in\R$).
Thus,
\[
x_{j}=\sum_{i=1}^{r}(j\lambda_{i}-[j\lambda_{i}])e_{i},
\]
from which it follows that $x_{j}\in P$ and, by \eqref{eq:IV.1.2},
$x_{j}\in P\cap H$. If one notices that
$x=x_{1}+\sum_{i=1}^{r}[\lambda_{i}]e_{i}$, one sees that the $\Z$-module $H$
is generated by $P\cap H$ and is thus of finite type.

On the other hand, since $P\cap H$ is finite and $\Z$ is infinite, there
exist distinct integers $j$ and $k$ such that $x_{j}=x_{k}$. It follows from
\eqref{eq:IV.1.2} that
$(j-k)\lambda_{i}=[j\lambda_{i}]-[k\lambda_{i}]$, which implies that the
$\lambda_{i}$'s are rational. Thus the $\Z$-module $H$ is generated by a
finite number of elements which are linear combinations with rational
coefficients of the $e_{i}$'s. Let $d$ be a common denominator
($d\in\Z-(0)$) of these coefficients. Clearly,
$dH\subset\sum_{i=1}^{r}\Z e_{i}$. Thus there exists a base $(f_{i})$ of the
$\Z$-module $\sum_{i=1}^{r}\Z e_{i}$ and integers $\alpha_{i}$ such that
$(\alpha_{1}f_{1},\ldots,\alpha_{r}f_{r})$ generates $dH$ (\hyperref[sec:1.5]{Chapter I, §5},
Theorem \ref{thm:I.5.1}). Since the $\Z$-module $dH$ has the same rank as $H$
and since $H\supset\sum_{i=1}^{r}\Z e_{i}$, the rank of $dH$ is $\geq r$.
Therefore, the rank of $dH$ equals $r$ and the $\alpha_{i}$'s are non-zero.
We may conclude that the $f_{i}$'s are, like the $e_{i}$'s, linearly
independent over $\R$. The module $dH$, and consequently $H$ itself, is
generated (over $\Z$) by $r$ vectors linearly independent over $\R$.
\end{proof}

An application of the preceding. Let
$t=(\theta_{1},\ldots,\theta_{n})\in\R^{n}$ such that at least one of the
$\theta_{i}$'s is irrational. Write $(e_{1},\ldots,e_{n})$ for the canonical
base of $\R^{n}$ and let $H$ denote the subgroup of $\R^{n}$ generated by
$(e_{1},\ldots,e_{n},t)$. The group $H$ is not discrete; otherwise the
methods employed in the proof of Theorem \ref{thm:IV.1.1} would provide us
with an expression for $t$ as a rational linear combination of the
$e_{i}$'s, an absurdity. Therefore, for any $\epsilon>0$ there exists a
non-zero element of $H$ whose distance from $0$ is smaller than $\epsilon$.
Therefore, there exist integers $p_{i}\in\Z$, $q\in\N$, $q\neq0$ such that
$|q\theta_{i}-p_{i}|\leq\epsilon$, which means that
\[
\Bigl|\theta_{i}-\frac{p_{i}}{q}\Bigr|\leq\frac{\epsilon}{q}
\quad\text{for all }i=1,\ldots,n.
\]
Let us remark that, simply by picking the multiple $n_{i}/q$ of $1/q$ nearest
$\theta_{i}$, we would obtain the cruder approximation
\[
\Bigl|\theta_{i}-\frac{n_{i}}{q}\Bigr|\leq\frac{1}{2q}
\quad(n_{i}\in\Z),\quad\text{for any }q>0.
\]

The result proved in the last paragraph is a basic theorem in the very rich
theory of approximation of irrational numbers by rationals. The reader
interested in learning more about this subject should consult Koksma,
``Diophantische Approximation'', Berlin (Springer), 1936.

\begin{definition}
\label{def:IV.1.1}
A discrete subgroup of rank $n$ of $\R^{n}$ is called a \emph{lattice} in
$\R^{n}$.
\end{definition}

By Theorem \ref{thm:IV.1.1} a lattice is generated over $\Z$ by a base of
$\R^{n}$, which is then a $\Z$-base for the given lattice. For each
$\Z$-base $e=(e_{1},\ldots,e_{n})$ of a lattice $H$ we shall write $P_{e}$
for the half open parallelotope
\[
P_{e}=\Bigl\{x\in\R^{n}\mid x=\sum_{i=1}^{n}\alpha_{i}e_{i},\quad
\text{with }0\leq\alpha_{i}<1\Bigr\}.
\]
Thus every point of $\R^{n}$ is congruent modulo $H$ to one and only one
point of $P_{e}$ for any fixed $e$ (we say, in this case, that $P_{e}$ is a
\emph{fundamental domain} for $H$). We shall write $\mu$ to denote the
\emph{Lebesgue measure} in $\R^{n}$, i.e.\ if $S$ is a measurable subset of
$\R^{n}$, $\mu(S)$ will stand for its measure (which we will also call its
volume).

\begin{lemma}
\label{lem:IV.1.1}
The volume $\mu(P_{e})$ is independent of the base $e$ chosen for $H$.
\end{lemma}

\begin{proof}
Let $f=(f_{1},\ldots,f_{n})$ be another base of $H$. Then
\[
f_{i}=\sum_{j=1}^{n}\alpha_{ij}e_{j}\quad\text{with }\alpha_{ij}\in\Z.
\]
By calculus we know that $\mu(P_{f})=|\det(\alpha_{ij})|\mu(P_{e})$. The
matrix $(\alpha_{ij})$, being associated with a change of base, is
invertible with an integer inverse matrix, so $\det(\alpha_{ij})=\pm1$.
Thus, $\mu(P_{f})=\mu(P_{e})$.
\end{proof}

The volume of the parallelotope $P_{e}$ associated with any base $e$ of $H$
is called the \emph{volume of the lattice} $H$ and is denoted $v(H)$ (the
word ``volume'' is here an abuse of language since $\mu(H)=0$, would it
perhaps be better to speak of the ``mesh'' of the lattice $H$?).

\begin{theorem}
\label{thm:IV.1.2}
\textnormal{(Minkowski)} Let $H$ be a lattice in $\R^{n}$ and let $S$ be a measurable subset of
$\R^{n}$ such that $\mu(S)>v(H)$. Then there exist two distinct points
$x,y\in S$ such that $x-y\in H$.
\end{theorem}

\begin{proof}
Let $e=(e_{1},\ldots,e_{n})$ be a $\Z$-base of $H$ and let $P_{e}$ be the
parallelotope associated with $e$. Since $P_{e}$ is a fundamental domain for
$H$, $S$ is the disjoint union of subsets of the form $S\cap(h+P_{e})$,
$h\in H$. It follows that
\begin{equation}
\mu(S)=\sum_{h\in H}\mu(S\cap(h+P_{e})).
\label{eq:IV.1.3}
\end{equation}
Since $\mu$ is translation invariant,
\[
\mu(S\cap(h+P_{e}))=\mu((-h+S)\cap P_{e}).
\]
The sets $(-h+S)\cap P_{e}$ $(h\in H)$ cannot all be pairwise disjoint,
otherwise
$\mu(P_{e})\geq\sum_{h\in H}\mu((-h+S)\cap P_{e})$, which contradicts
\eqref{eq:IV.1.3} and the hypothesis $\mu(P_{e})=v(H)<\mu(S)$. Consequently,
there exist two distinct elements $h$ and $h'$ of $H$ such that
$P_{e}\cap(-h+S)\cap(-h'+S)\neq\emptyset$. Let $x$ and $y$ be elements of
$S$ such that $-h+x=-h'+y$. Then $x-y=h-h'\in H$ and $x\neq y$, since
$h\neq h'$.
\end{proof}

\begin{statement}{Corollary}
Let $H$ be a lattice in $\R^{n}$ and let $S$ be a measurable subset of
$\R^{n}$ which is symmetric with respect to $0$ and convex. Assume that $S$
satisfies at least one of the following two conditions:
\begin{alist}
\item $\mu(S)>2^{n}v(H)$
\item $\mu(S)\geq2^{n}v(H)$ and $S$ is compact.
\end{alist}
Then $S\cap(H-\{0\})\neq\emptyset$.
\end{statement}

\begin{proof}
In case (a) apply Theorem \ref{thm:IV.1.2} to the set $S'=\frac{1}{2}S$
($\mu(S')=2^{-n}\mu(S)>v(H)$). Let $y$ and $z$ be distinct points of $S'$
such that $y-z\in H$. Then $y-z$ also belongs to $S$, because
$y-z=\frac{1}{2}(2y+(-2z))$ and $S$ is both symmetric and convex. Therefore,
$y-z\in S\cap(H-\{0\})$. To prove that case (b) also implies the conclusion
of the corollary apply case (a) to $(1+\epsilon)S$ for $\epsilon>0$. Thus,
$(H-\{0\})\cap(1+\epsilon)S$ is a non-empty compact set (it is even finite,
since it is compact and discrete). Furthermore,
$\bigcap_{\epsilon>0}(H-\{0\})\cap(1+\epsilon)S\neq\emptyset$, since an
intersection of non-empty, nested compact sets is never void. This means
that there is a point of $H-\{0\}$ which belongs to $(1+\epsilon)S$ for all
$\epsilon>0$, therefore, since $S$ is compact, it belongs to $S$, too.
\end{proof}

The hypothesis of compactness is needed in (b). Consider, as a
counter-example to (b) with compactness omitted, the open parallelotope
\[
\Bigl\{x\in\R^{n}\mid x=\sum_{i=1}^{n}\lambda_{i}e_{i},\quad
-1<\lambda_{i}<1\Bigr\}
\]
built on the base $e=(e_{1},\ldots,e_{n})$, and the lattice having $e$ as a
$\Z$-base.

%% ----------------------------------------------------------------------
\sechead{2}{The canonical imbedding of a number field}%
        {THE CANONICAL IMBEDDING OF A NUMBER FIELD}

Let $K$ be a number field and let $n$ be its degree. We have seen
(\hyperref[sec:2.4]{Chapter II, §4}, Theorem \ref{thm:II.4.1}) that there are $n$ distinct
isomorphisms $\sigma_{i}:K\to\C$. There are exactly $n$, because the
minimal polynomial for a primitive element of $K$ over $\Q$ (ibid.,
corollary of Theorem \ref{thm:II.4.1}) has only $n$ roots in $\C$. Let
$\alpha:\C\to\C$ be complex conjugation. Then, for any $i=1,\ldots,n$,
$\alpha\circ\sigma_{i}=\sigma_{j}$, $1\leq j\leq n$, and
$\sigma_{i}=\sigma_{j}$ if and only if $\sigma_{i}(K)\subset\R$. Write
$r_{1}$ for the number of indices such that $\sigma_{i}(K)\subset\R$. Then
$n-r_{1}$ is an even number, so we may write
\begin{equation}
r_{1}+2r_{2}=n.
\label{eq:IV.2.1}
\end{equation}
Let us renumber the $\sigma_{i}$'s so that $\sigma_{i}(K)\subset\R$ for
$1\leq i\leq r_{1}$ and so that
$\sigma_{j+r_{2}}(x)=\overline{\sigma_{j}(x)}$ for
$r_{1}+1\leq j\leq r_{1}+r_{2}$. Then the first $r_{1}+r_{2}$ isomorphisms
determine the last $r_{2}$. For $x\in K$ we define
\begin{equation}
\sigma(x)=\bigl(\sigma_{1}(x),\ldots,
\sigma_{r_{1}+r_{2}}(x)\bigr)\in\R^{r_{1}}\times\C^{r_{2}}.
\label{eq:IV.2.2}
\end{equation}
We call $\sigma$ the \emph{canonical imbedding} of $K$ in
$\R^{r_{1}}\times\C^{r_{2}}$; it is an injective ring homomorphism. We shall
frequently identify $\R^{r_{1}}\times\C^{r_{2}}$ with $\R^{n}$
(cf.\ \eqref{eq:IV.2.1}). The notations $\sigma$, $K$, $n$, $r_{1}$, and
$r_{2}$ will be in use for the rest of this chapter.

\begin{proposition}
\label{prop:IV.2.1}
If $M$ is a free $\Z$-submodule of $K$ of rank $n$ and if
$(x_{i})_{1\leq i\leq n}$ is a $\Z$-base of $M$, then $\sigma(M)$ is a
lattice in $\R^{n}$, whose volume is
\begin{equation}
v(\sigma(M))=2^{-r_{2}}
\Bigl|\det_{1\leq i,j\leq n}\bigl(\sigma_{i}(x_{j})\bigr)\Bigr|.
\label{eq:IV.2.3}
\end{equation}
\end{proposition}

\begin{proof}
For fixed $i$ the coordinates of $\sigma(x_{i})$ with respect to the
canonical base of $\R^{n}$ are
\begin{equation}
\begin{gathered}
\sigma_{1}(x_{i}),\ldots,\sigma_{r_{1}}(x_{i}),
\mathrm{R}(\sigma_{r_{1}+1}(x_{i})),\mathrm{I}(\sigma_{r_{1}+1}(x_{i})),\ldots,\\
\mathrm{R}(\sigma_{r_{1}+r_{2}}(x_{i})),\mathrm{I}(\sigma_{r_{1}+r_{2}}(x_{i})),
\end{gathered}
\label{eq:IV.2.4}
\end{equation}
where $\mathrm{R}$ and $\mathrm{I}$ denote, respectively, the real and
imaginary parts. Let us calculate the determinant $D$ of the matrix whose
$i$th column is given by \eqref{eq:IV.2.4}. Making use of the formulas
$\mathrm{R}(z)=\frac{1}{2}(z+\overline{z})$ and
$\mathrm{I}(z)=(1/2i)(z-\overline{z})$ for $z\in\C$ and of the linearity
of both $\mathrm{R}$ and $\mathrm{I}$, we obtain
$D=(2i)^{-r_{2}}\det(\sigma_{j}(x_{i}))$. Since the
$x_{i}$'s form a base for $K$ over $\Q$, $\det(\sigma_{j}(x_{i}))\neq0$
(\hyperref[sec:2.7]{Chapter II, §7}, Proposition \ref{prop:II.7.3}) and therefore $D\neq0$. Thus
the vectors $\sigma(x_{i})$ are linearly independent in $\R^{n}$, so that
the $\Z$-module which they generate (call it $\sigma(M)$) is a lattice in
$\R^{n}$. The calculation of $D$ given above shows that \eqref{eq:IV.2.3}
does give its volume.
\end{proof}

\begin{proposition}
\label{prop:IV.2.2}
Let $d$ be the absolute discriminant of $K$, let $A$ be the ring of integers
in $K$, and let $\mathfrak{a}$ be a non-zero integral ideal of $A$. Then
$\sigma(A)$ and $\sigma(\mathfrak{a})$ are lattices. Moreover,
\begin{equation}
v(\sigma(A))=2^{-r_{2}}|d|^{1/2}\quad\text{and}\quad
v(\sigma(\mathfrak{a}))=2^{-r_{2}}|d|^{1/2}\mathrm{N}(\mathfrak{a}).
\label{eq:IV.2.5}
\end{equation}
\end{proposition}

\begin{proof}
We know that $A$ and $\mathfrak{a}$ are free $\Z$-modules of rank $n$, so we
may apply Proposition \ref{prop:IV.2.1}. On the other hand, if $(x_{i})$ is
a $\Z$-base for $A$, then $d=\det(\sigma_{i}(x_{j}))^{2}$ (\hyperref[sec:2.7]{Chapter II, §7},
Proposition \ref{prop:II.7.3}). This proves the first formula in
\eqref{eq:IV.2.5}. The second formula follows from the first and the
observation that $\sigma(\mathfrak{a})$ is a subgroup of $\sigma(A)$ of
index $\mathrm{N}(\mathfrak{a})$ (\hyperref[sec:3.5]{Chapter III, §5}, \hyperref[def:III.5.1]{Definition 1}). A fundamental
domain for $\sigma(\mathfrak{a})$ may obviously be constructed as the
disjoint union of $\mathrm{N}(\mathfrak{a})$ copies of a fundamental domain for
$\sigma(A)$.
\end{proof}

%% ----------------------------------------------------------------------
\sechead{3}{Finiteness of the ideal class group}%
        {FINITENESS OF THE IDEAL CLASS GROUP}

\begin{proposition}
\label{prop:IV.3.1}
Let $K$ be a number field, $n$ its degree, $r_{1}$ and $r_{2}$ the integers
defined in the beginning of \hyperref[sec:4.2]{§2}, $d$ the absolute discriminant of $K$, and
$\mathfrak{a}$ a non-zero integral ideal of $K$. Then $\mathfrak{a}$ contains
a non-zero element $x$ such that
\begin{equation}
|\mathrm{N}_{K/\Q}(x)|\leq\Bigl(\frac{4}{\pi}\Bigr)^{r_{2}}\frac{n!}{n^{n}}
|d|^{1/2}\mathrm{N}(\mathfrak{a}).
\label{eq:IV.3.1}
\end{equation}
\end{proposition}

\begin{proof}
Let $\sigma$ be the canonical imbedding of $K$ into
$\R^{r_{1}}\times\C^{r_{2}}$ (\hyperref[sec:4.2]{§2}). Let $t$ be a positive real number and let
$B_{t}$ be the set of all elements
$(y_{1},\ldots,y_{r_{1}},z_{1},\ldots,z_{r_{2}})\in
\R^{r_{1}}\times\C^{r_{2}}$ such that
\begin{equation}
\sum_{i=1}^{r_{1}}|y_{i}|+2\sum_{j=1}^{r_{2}}|z_{j}|\leq t.
\label{eq:IV.3.2}
\end{equation}
Then $B_{t}$ is a set which is compact and convex and symmetric relative to
$0\in\R^{n}$. By a calculation given in the appendix to this chapter
\begin{equation}
\mu(B_{t})=2^{r_{1}}\Bigl(\frac{\pi}{2}\Bigr)^{r_{2}}\frac{t^{n}}{n!}.
\label{eq:IV.3.3}
\end{equation}
Now choose $t$ such that
$\mu(B_{t})=2^{n}v(\sigma(\mathfrak{a}))$, i.e.\ such that
\[
2^{r_{1}}\Bigl(\frac{\pi}{2}\Bigr)^{r_{2}}\frac{t^{n}}{n!}
=2^{n-r_{2}}|d|^{1/2}\mathrm{N}(\mathfrak{a})
\quad(\text{\S2, Proposition \ref{prop:IV.2.2}}),
\]
or such that
$t^{n}=2^{n-r_{1}}\pi^{-r_{2}}n!\,|d|^{1/2}\mathrm{N}(\mathfrak{a})$.

By the corollary to Theorem \ref{thm:IV.1.2}, \hyperref[sec:4.1]{§1}, there exists a non-zero
element $x\in\mathfrak{a}$ such that $\sigma(x)\in B_{t}$. Its norm has
absolute value
\[
|\mathrm{N}(x)|=\prod_{i=1}^{r_{1}}|\sigma_{i}(x)|
\prod_{j=r_{1}+1}^{r_{1}+r_{2}}|\sigma_{j}(x)|^{2}.
\]
By the fact that the geometric mean never exceeds the arithmetic mean
(cf.\ Hardy-Littlewood-Polya, ``Inequalities'', Cambridge University Press)
we have
\[
|\mathrm{N}(x)|\leq\Bigl[\frac{1}{n}\sum_{i=1}^{r_{1}}|\sigma_{i}(x)|
+\frac{2}{n}\sum_{j=r_{1}+1}^{r_{1}+r_{2}}|\sigma_{j}(x)|\Bigr]^{n}
\leq\frac{t^{n}}{n^{n}}\quad(\text{by \eqref{eq:IV.3.2}}).
\]
Consequently,
\[
|\mathrm{N}(x)|\leq\frac{1}{n^{n}}2^{n-r_{1}}\pi^{-r_{2}}n!\,|d|^{1/2}
\mathrm{N}(\mathfrak{a}),
\]
which, combined with the relation $r_{1}+2r_{2}=n$, yields
\eqref{eq:IV.3.1}.
\end{proof}

\begin{corollary}
\label{cor:IV.3.1}
With the same notations, every ideal class of $K$ (\hyperref[sec:3.4]{Chapter III, §4})
contains an integral ideal $\mathfrak{b}$ such that
\begin{equation}
\mathrm{N}(\mathfrak{b})\leq\Bigl(\frac{4}{\pi}\Bigr)^{r_{2}}\frac{n!}{n^{n}}
|d|^{1/2}.
\label{eq:IV.3.4}
\end{equation}
\end{corollary}

\begin{proof}
Let $\mathfrak{a}'$ be an ideal of the given class. By multiplying
$\mathfrak{a}'$ by a principal ideal we may suppose that
$\mathfrak{a}=\mathfrak{a}'^{-1}$ is an integral ideal. Take a non-zero
element $x\in\mathfrak{a}$ for which \eqref{eq:IV.3.1} holds. Then
$\mathfrak{b}=x\mathfrak{a}^{-1}$ is an integral ideal in the same class as
$\mathfrak{a}'$, and $\mathrm{N}(\mathfrak{b})$ satisfies \eqref{eq:IV.3.4} by virtue
of the multiplicativity of norms (\hyperref[sec:3.5]{Chapter III, §5}, Proposition
\ref{prop:III.5.2}).
\end{proof}

\begin{corollary}
\label{cor:IV.3.2}
Let $K$ be a number field, let $n$ be its degree, and let $d$ be its
absolute discriminant. Then, for $n\geq2$,
\[
|d|\geq\frac{\pi}{3}\Bigl(\frac{3\pi}{4}\Bigr)^{n-1}
\]
and $n/(\log|d|)$ is majorized by a constant independent of $K$.
\end{corollary}

Since there is always a non-zero integral ideal $\mathfrak{b}$ in $K$ and
$\mathrm{N}(\mathfrak{b})\geq1$, we obtain, from \eqref{eq:IV.3.4},
\[
|d|^{1/2}\geq\Bigl(\frac{\pi}{4}\Bigr)^{r_{2}}\frac{n^{n}}{n!}.
\]
From $\pi/4<1$ and $2r_{2}\leq n$ we conclude that $|d|\geq a_{n}$, where
$a_{n}=(\pi/4)^{n}[n^{2n}/(n!)^{2}]$. We observe that
\[
a_{2}=\frac{\pi^{2}}{4}\quad\text{and}\quad
\frac{a_{n+1}}{a_{n}}=\frac{\pi}{4}\Bigl(1+\frac{1}{n}\Bigr)^{2n}
=\frac{\pi}{4}(1+2+\text{positive terms})
\]
(by use of the binomial formula), so $a_{n+1}/a_{n}\geq3\pi/4$. Hence, for
$n\geq2$,
\[
|d|\geq\frac{\pi^{2}}{4}\Bigl(\frac{3\pi}{4}\Bigr)^{n-2},
\]
from which the inequality statement in the corollary is immediate. The
uniform majoration of $n/(\log|d|)$ follows by taking logarithms.

\begin{theorem}
\label{thm:IV.3.1}
\textnormal{(Hermite-Minkowski)} For any number field $K\neq\Q$, the
absolute discriminant $d$ of $K$ is $\neq\pm1$.
\end{theorem}

\begin{proof}
Using Corollary \ref{cor:IV.3.2} we see that
$|d|\geq(\pi/3)(3\pi/4)^{n-1}$. Since $\pi/3>1$ and $3\pi/4>1$, we have
$|d|>1$.
\end{proof}

\begin{theorem}
\label{thm:IV.3.2}
\textnormal{(Dirichlet)} For any number field $K$ the ideal class group is
finite (\hyperref[sec:3.4]{Chapter III, §4}).
\end{theorem}

\begin{proof}
By Corollary \ref{cor:IV.3.1} of Proposition \ref{prop:IV.3.1} it suffices
to show that, for every positive integer $q$, the set of all integral ideals
$\mathfrak{b}$ of $K$ which have $q$ as their norms is a finite set. For
such an ideal $\mathfrak{b}$ we have
$\operatorname{card}(A/\mathfrak{b})=q$ (\hyperref[sec:3.5]{Chapter III, §5}). It follows that
$q\in\mathfrak{b}$, since for any group the order of an element divides the
order of the group. Thus our ideals $\mathfrak{b}$ are among those which
contain $Aq$, and there can be only finitely many such ideals (\hyperref[sec:3.4]{Chapter III,
§4}, formula \eqref{eq:III.4.8}; or the finiteness of $A/Aq$).
\end{proof}

\begin{theorem}
\label{thm:IV.3.3}
\textnormal{(Hermite)} In $\C$ there are only finitely many number fields
with a given discriminant $d$.
\end{theorem}

\begin{proof}
By Corollary \ref{cor:IV.3.2} of Proposition \ref{prop:IV.3.1} the degree of
such a field is bounded. We may suppose $n$ and the integers $r_{1}$ and
$r_{2}$ are fixed. Let $K$ be such a field.

In $\R^{r_{1}}\times\C^{r_{2}}$ consider the following set $B$:

(a) If $r_{1}>0$, $B$ is the set of all elements
$(y_{1},\ldots,y_{r_{1}},z_{1},\ldots,z_{r_{2}})\in
\R^{r_{1}}\times\C^{r_{2}}$ such that
\begin{equation}
|y_{1}|\leq2^{n}\Bigl(\frac{\pi}{2}\Bigr)^{-r_{2}}|d|^{1/2};
\label{eq:IV.3.5}
\end{equation}
$|y_{i}|\leq\frac{1}{2}$ for $i=2,\ldots,r_{1}$; and
$|z_{j}|\leq\frac{1}{2}$ for $j=1,\ldots,r_{2}$.

(b) If $r_{1}=0$, $B$ is the set of all elements
$(z_{1},\ldots,z_{r_{2}})\in\C^{r_{2}}$ such that
\begin{equation}
|z_{1}-\overline{z}_{1}|\leq2^{n}\Bigl(\frac{\pi}{2}\Bigr)^{1-r_{2}}
|d|^{1/2},\quad |z_{1}+\overline{z}_{1}|\leq\frac{1}{2},\quad
\text{and}\quad |z_{j}|\leq\frac{1}{2}\ \text{for }j=2,\ldots,r_{2}.
\label{eq:IV.3.6}
\end{equation}
Clearly, $B$ is a compact and convex set, which is symmetric about the
origin in $\R^{n}$ and which has volume $2^{n-r_{2}}|d|^{1/2}$. Writing
$\sigma$ for the canonical imbedding of $K$ (\hyperref[sec:4.2]{§2}), we obtain by means of
Proposition \ref{prop:IV.2.2} of \hyperref[sec:4.2]{§2} and the corollary to Theorem
\ref{thm:IV.1.2} of \hyperref[sec:4.1]{§1} an integer $x\neq0$ of $K$ for which $\sigma(x)\in B$.

Let us show that $x$ is a primitive element of $K$ over $\Q$. In case (a),
\eqref{eq:IV.3.5} shows that $|\sigma_{i}(x)|\leq\frac{1}{2}$ for $i\neq1$.
Since $|\mathrm{N}(x)|=\prod_{i=1}^{n}|\sigma_{i}(x)|$ is a positive integer
(\hyperref[sec:2.6]{Chapter II, §6}, corollary of Proposition \ref{prop:II.6.2}), we may
conclude that $|\sigma_{1}(x)|\geq1$, so $\sigma_{1}(x)\neq\sigma_{i}(x)$
for all $i\neq1$. However, if $x$ were not primitive, $\sigma_{1}(x)$ would
coincide with $\sigma_{i}(x)$ for some $i\neq1$ (\hyperref[sec:2.6]{Chapter II, §6},
Proposition \ref{prop:II.6.1}). In case (b) we see similarly that
$|\sigma_{1}(x)|=|\overline{\sigma_{1}(x)}|\geq1$, so
$\sigma_{1}(x)\neq\sigma_{j}(x)$ when $\sigma_{j}$ is not $\sigma_{1}$ or
$\overline{\sigma_{1}}$. Now \eqref{eq:IV.3.6} implies that the real part
$\mathrm{R}(\sigma_{1}(x))\leq\frac{1}{2}$. But this means that $\sigma_{1}(x)$
cannot be real, so $\sigma_{1}(x)\neq\overline{\sigma_{1}(x)}$. As in case
(a) we conclude that $x$ is primitive.

Formulas \eqref{eq:IV.3.5} and \eqref{eq:IV.3.6} imply that the conjugates
$\sigma_{i}(x)$ of $x$ are bounded. Therefore the elementary symmetric
functions of the $\sigma_{i}(x)$'s are also bounded. In other words, the
coefficients, as well as the degree, of the minimal polynomial of $x$ are
bounded. Since $x$ is an integer, its minimal polynomial is a monic
polynomial with coefficients in $\Z$ (\hyperref[sec:2.6]{Chapter II, §6}, the corollary of
Proposition \ref{prop:II.6.2}). The degree and the coefficients of the
minimal polynomial of $x$ being bounded, there are only finitely many
possibilities for the minimal polynomial of $x$, consequently only finitely
many possible values for $x\in\C$. As $x$ generates $K$, there are only
finitely many possibilities for $K$.
\end{proof}

%% ----------------------------------------------------------------------
\sechead{4}{The unit theorem}{THE UNIT THEOREM}

Let $K$ be a number field and let $A$ be the ring of integers in $K$. By
abuse of language we use the expression ``units of $K$'' to refer to the
units in the ring $A$. We remind the reader that in any ring the units form
a group under multiplication. We write $A^{*}$ for the group of units in
$A$.

The following result will be useful.

\begin{proposition}
\label{prop:IV.4.1}
Let $K$ be a number field and let $x\in K$. In order that $x$ be a unit of
$K$ it is necessary and sufficient that $x$ be an integer of $K$ of norm
$\pm1$.
\end{proposition}

\begin{proof}
If $x$ is a unit of $K$, then $\mathrm{N}(x)$ and $\mathrm{N}(x^{-1})$
belong to $\Z$. We have
$\mathrm{N}(x)\mathrm{N}(x^{-1})=\mathrm{N}(xx^{-1})=1$, so
$\mathrm{N}(x)=\pm1$. Conversely, let $x$ be an integer of $K$ with norm
$\pm1$. Its characteristic equation has the form
\[
x^{n}+a_{n-1}x^{n-1}+\cdots+a_{1}x\pm1=0
\quad\text{with }a_{i}\in\Z\ \text{(\hyperref[sec:2.6]{Chapter II, §6})}.
\]
Thus, $\pm(x^{n-1}+a_{n-1}x^{n-2}+\cdots+a_{1})=x^{-1}$ and, since
$x^{-1}$ is an integer of $K$, $x$ is a unit.
\end{proof}

\begin{theorem}
\label{thm:IV.4.1}
\textnormal{(Dirichlet)} Let $K$ be a number field, $n$ its degree, and let
$r_{1}$ and $r_{2}$ be the integers defined in \hyperref[sec:4.2]{§2}. Set
$r=r_{1}+r_{2}-1$. The group $A^{*}$ of units of $K$ is isomorphic to
$\Z^{r}\times G$, where $G$ is a finite cyclic group comprised of the roots
of unity contained in $K$.
\end{theorem}

\begin{proof}
First we shall show that $A^{*}$ is a commutative group of finite type. Then
we shall calculate its rank.

Consider the canonical imbedding (\hyperref[sec:4.2]{§2})
$x\mapsto(\sigma_{1}(x),\ldots,\sigma_{r_{1}+r_{2}}(x))$ of $K$ into
$\R^{r_{1}}\times\C^{r_{2}}$ and the mapping
\begin{equation}
x\mapsto L(x)=(\log|\sigma_{1}(x)|,\ldots,
\log|\sigma_{r_{1}+r_{2}}(x)|)
\label{eq:IV.4.1}
\end{equation}
of $K^{*}$ to $\R^{r_{1}+r_{2}}$. \eqref{eq:IV.4.1} is a homomorphism
(i.e.\ $L(xy)=L(x)+L(y)$), called the \emph{logarithmic imbedding} of
$K^{*}$. Let $B$ be a compact subset of $\R^{r_{1}+r_{2}}$. Let us show that
the set $B'$ of units $x\in A^{*}$ such that $L(x)\in B$ is a finite set.
Indeed, since $B$ is bounded, there exists a real number $\alpha>1$ such
that, for all $x\in B'$, $\alpha^{-1}\leq|\sigma_{i}(x)|\leq\alpha$
($i=1,\ldots,n$). It follows that the elementary symmetric functions of the
$\sigma_{i}(x)$'s are bounded in absolute value. Since they belong to $\Z$
(because $x\in A$), the set of possible values for the symmetric functions
of the $\sigma_{i}(x)$'s is a finite set. Therefore, there are only finitely
many possible characteristic polynomials for elements $x\in B'$ and
consequently only finitely many possible values for $x$. Thus, $B'$ is a
finite set.

The finiteness of $B'$ implies, immediately, the following statements:

(a) The kernel $G$ of the restriction of $L$ to $A^{*}$ is a finite group.
It consists, therefore, of roots of unity and is cyclic (\hyperref[sec:1.6]{Chapter I, §6},
Theorem \ref{thm:I.6.1}). Clearly, every root of unity in $K$ belongs to the
kernel of $L$, for the roots of unity in $K$ are integers and, if $x$ is a
root of unity in $K$,
$|\sigma_{i}(x)|^{q}=|\sigma_{i}(x^{q})|=|1|=1$, so
$|\sigma_{i}(x)|=1$ for any $i$.

(b) The image $L(A^{*})$ is a discrete subgroup of $\R^{r_{1}+r_{2}}$ (\hyperref[sec:4.1]{§1}).
Consequently, $L(A^{*})$ is a free $\Z$-module of rank
$s\leq r_{1}+r_{2}$ (\hyperref[sec:4.1]{§1}, Theorem \ref{thm:IV.1.1}). Since $L(A^{*})$ is
free, $A^{*}$ is isomorphic to $G\times L(A^{*})=G\times\Z^{s}$. It remains
to show that the rank $s$ of $L(A^{*})$ equals $r_{1}+r_{2}-1$.

The inequality $s\leq r_{1}+r_{2}-1$ is easy. Indeed, for $x\in A^{*}$, the
relation
\[
\pm1=\mathrm{N}(x)=\prod_{i=1}^{n}\sigma_{i}(x)
=\prod_{i=1}^{r_{1}}\sigma_{i}(x)
\prod_{j=r_{1}+1}^{r_{1}+r_{2}}\sigma_{j}(x)\overline{\sigma_{j}(x)}
\quad(\text{Proposition \ref{prop:IV.4.1}})
\]
implies that the vector $L(x)=(y_{1},\ldots,y_{r_{1}+r_{2}})$ lies in the
hyperplane $W$ defined by the equation
\begin{equation}
\sum_{i=1}^{r_{1}}y_{i}+2\sum_{j=r_{1}+1}^{r_{1}+r_{2}}y_{j}=0.
\label{eq:IV.4.2}
\end{equation}
Since $L(A^{*})$ is a discrete subgroup of $W$, $s\leq r_{1}+r_{2}-1$.

Now we show that $L(A^{*})$ contains $r=r_{1}+r_{2}-1$ linearly independent
vectors. This requires a more delicate argument. We are going to show that
for any linear form $f\neq0$ on $W$, there exists a unit $u$ such that
$f(L(u))\neq0$. As the projection of $W$ on $\R^{r}$ is an isomorphism (by
\eqref{eq:IV.4.2}), we may write, for any
$y=(y_{1},\ldots,y_{r})\in W\subset\R^{r_{1}+r_{2}}$,
\begin{equation}
f(y)=c_{1}y_{1}+\cdots+c_{r}y_{r}\quad\text{with }c_{i}\in\R.
\label{eq:IV.4.3}
\end{equation}
Fix a real number $\alpha$ such that
$\alpha\geq2^{n}(1/2\pi)^{r_{2}}|d|^{1/2}$. For any set
$\lambda=(\lambda_{1},\ldots,\lambda_{r})$ of $r$ positive real numbers take
$\lambda_{r+1}>0$ such that
$\prod_{i=1}^{r_{1}}\lambda_{i}
\prod_{j=r_{1}+1}^{r_{1}+r_{2}}\lambda_{j}^{2}=\alpha$. In
$\R^{r_{1}}\times\C^{r_{2}}$ the set $B$ of elements
$(y_{1},\ldots,y_{r_{1}},z_{1},\ldots,z_{r_{2}})$ ($y_{i}\in\R$ and
$z_{j}\in\C$) such that $|y_{i}|\leq\lambda_{i}$ and
$|z_{j}|\leq\lambda_{j+r_{1}}$ is compact, convex, symmetric about $0$, and
of volume
$\prod_{i=1}^{r_{1}}2\lambda_{i}
\prod_{j=1+r_{1}}^{r_{2}+r_{1}}\pi\lambda_{j}^{2}
=2^{r_{1}}\pi^{r_{2}}\alpha\geq2^{r_{1}+n-r_{2}}|d|^{1/2}$.
\footnote{One may calculate this volume by making use of the observation
that $B$ is a product of intervals, of discs, and of a rectangle in case
(b).} It follows
from Proposition \ref{prop:IV.2.2} of \hyperref[sec:4.2]{§2} and the corollary to Theorem
\ref{thm:IV.1.2} of \hyperref[sec:4.1]{§1} that there exists an integer $x_{\lambda}$ of $K$
such that $\sigma(x_{\lambda})\in B$. This means that
$|\sigma_{i}(x_{\lambda})|\leq\lambda_{i}$ for $i=1,\ldots,n$ (putting
$\lambda_{j+r_{2}}=\lambda_{j}$ for $j=r_{1}+1,\ldots,r_{1}+r_{2}$). Since
$x_{\lambda}$ is an integer,
\[
1\leq|\mathrm{N}(x_{\lambda})|=\prod_{i=1}^{n}|\sigma_{i}(x_{\lambda})|
\leq\prod_{i=1}^{r_{1}}\lambda_{i}
\prod_{j=r_{1}+1}^{r_{1}+r_{2}}\lambda_{j}^{2}=\alpha.
\]
On the other hand, for any $i$
\[
|\sigma_{i}(x_{\lambda})|=|\mathrm{N}(x_{\lambda})|
\prod_{j\neq i}|\sigma_{j}(x_{\lambda})|^{-1}
\geq\prod_{j\neq i}\lambda_{j}^{-1}=\lambda_{i}\alpha^{-1}.
\]
Now we have $\lambda_{i}\alpha^{-1}\leq|\sigma_{i}(x_{\lambda})|\leq
\lambda_{i}$ for any $i$, so that
\begin{equation}
0\leq\log\lambda_{i}-\log|\sigma_{i}(x_{\lambda})|\leq\log\alpha.
\label{eq:IV.4.4}
\end{equation}
Use of \eqref{eq:IV.4.3} entails
\begin{equation}
\Bigl|f(L(x_{\lambda}))-\sum_{i=1}^{r}c_{i}\log\lambda_{i}\Bigr|
\leq\Bigl(\sum_{i=1}^{r}|c_{i}|\Bigr)\log\alpha.
\label{eq:IV.4.5}
\end{equation}
Let $\beta$ be a constant which is strictly larger than the right-hand side
of \eqref{eq:IV.4.5} and, for every positive integer $h$, select $r$
positive real numbers $\lambda_{i,h}$ ($i=1,\ldots,r$) satisfying
$\sum_{i=1}^{r}c_{i}\log\lambda_{i,h}=2\beta h$. Put
$\lambda(h)=(\lambda_{1,h},\ldots,\lambda_{r,h})$ and let $x_{h}$ be the
corresponding integer $x_{\lambda(h)}$. By \eqref{eq:IV.4.5} we have
$|f(L(x_{h}))-2\beta h|<\beta$, so
\begin{equation}
(2h-1)\beta<f(L(x_{h}))<(2h+1)\beta.
\label{eq:IV.4.6}
\end{equation}
It follows from \eqref{eq:IV.4.6} that the numbers $f(L(x_{h}))$ are all
distinct. On the other hand, since $|\mathrm{N}(x_{h})|\leq\alpha$, there are only
finitely many distinct ideals of the form $Ax_{h}$ (cf.\ \hyperref[sec:4.3]{§3}, the proof of
Theorem \ref{thm:IV.3.2}). Therefore, there exist two distinct indices $h$
and $k$ such that $Ax_{h}=Ax_{k}$ and, consequently, a unit $u\in A$ such
that $x_{k}=ux_{h}$. We may conclude (since $f$ is linear) that
$f(L(u))=f(L(x_{k}))-f(L(x_{h}))\neq0$.
\end{proof}

\begin{statement}{Remark}
\upshape
Theorem \ref{thm:IV.4.1} (called the ``unit theorem'') implies that there
exist $r(=r_{1}+r_{2}-1)$ units $(u_{i})$ of $K$ such that any unit $u$ of
$K$ may be uniquely expressed in the form
\begin{equation}
u=zu_{1}^{n_{1}}\cdots u_{r}^{n_{r}}
\label{eq:IV.4.7}
\end{equation}
with $n_{i}\in\Z$ and $z$ a root of unity. The set $(u_{i})$,
$i=1,\ldots,r$, is called a \emph{fundamental system of units} of $K$.
\end{statement}

\begin{statement}{Example: cyclotomic fields}
\upshape
Let $p$ be an odd prime number, let $z$ be a primitive complex $p$th root of
unity, and let $K$ be the cyclotomic field $\Q[z]$ (cf.\ \hyperref[sec:2.9]{Chapter II, §9}). We
know that $[K:\Q]=p-1$ (ibid., Theorem \ref{thm:II.9.1}). Since no conjugate
of $z$ in $\C$ is real, $r_{1}=0$ and $2r_{2}=p-1$, so $r=(p-3)/2$.
\end{statement}

%% ----------------------------------------------------------------------
\sechead{5}{Units in imaginary quadratic fields}%
        {UNITS IN IMAGINARY QUADRATIC FIELDS}

Let $K$ be an imaginary quadratic field (\hyperref[sec:2.5]{Chapter II, §5}). Then $r_{1}=0$,
$2r_{2}=2$, $r_{2}=1$, and $r_{1}+r_{2}-1=0$. Thus the only units in $K$ are
the roots of unity contained in $K$ (\hyperref[sec:4.4]{§4}, Theorem \ref{thm:IV.4.1}), a finite
cyclic group. With a little calculation we shall prove this result directly
and make it more precise.

Let $K=\Q[\sqrt{-m}]$, where $m$ is a square-free positive integer. Recall
that the units of $K$ are integers of norm $\pm1$ (\hyperref[sec:4.4]{§4}, Proposition
\ref{prop:IV.4.1}).

(1) If $m\equiv1$ or $2\pmod{4}$, the ring of integers of $K$ is
$\Z+\Z\sqrt{-m}$ (\hyperref[sec:2.5]{Chapter II, §5}, Theorem \ref{thm:II.5.1}). For
$x=a+b\sqrt{-m}$ ($a,b\in\Z$) we have $\mathrm{N}(x)=a^{2}+mb^{2}\geq0$. In order
that $x$ be a unit we must have $a^{2}+mb^{2}=1$. If $m\geq2$, this implies
that $b=0$ and $a=\pm1$, so $x=\pm1$. If $m=1$, besides the solution
$x=\pm1$, there are the solutions $a=0$, $b=\pm1$, i.e.\ $x=\pm i$ (with
$i^{2}=-1$).

(2) If $m\equiv3\pmod{4}$, the ring of integers of $K$ is
$\Z+\Z[(1+\sqrt{-m})/2]$ (\hyperref[sec:2.5]{Chapter II, §5}, Theorem \ref{thm:II.5.1}). For
$x=a+(b/2)(1+\sqrt{-m})$ ($a,b\in\Z$), we have
$\mathrm{N}(x)=(a+b/2)^{2}+mb^{2}/4$. In order that $x$ be a unit we must have
$(2a+b)^{2}+mb^{2}=4$. If $m\geq7$, this implies that $b=0$, so $2a=\pm2$,
$a=\pm1$, and $x=\pm1$. If $m=3$, the relations $b=\pm1$ and
$(2a\pm1)^{2}=1$ entail the additional solutions
$x=\frac{1}{2}(\pm1\pm\sqrt{-3})$ (the signs $\pm$ being independent).

Summarizing, we have proved the following result:

\begin{proposition}
\label{prop:IV.5.1}
If $K$ is a quadratic imaginary field, the group $G$ of units in $K$ is
comprised of the square roots of unity, $+1$ and $-1$, except in the
following two cases:
\begin{enumerate}[label=(\arabic*),leftmargin=2.0em,itemsep=0.1ex,topsep=0.3ex]
\item If $K=\Q[i]$ ($i^{2}=-1$), $G$ is comprised of the fourth roots of
unity: $i,-1,-i,1$.
\item If $K=\Q[\sqrt{-3}]$, $G$ is comprised of the sixth roots of unity:
$[(1+\sqrt{-3})/2]^{j}$, $j=0,1,\ldots,5$.
\end{enumerate}
\end{proposition}

%% ----------------------------------------------------------------------
\sechead{6}{Units in real quadratic fields}%
        {UNITS IN REAL QUADRATIC FIELDS}

This section is going to be considerably more interesting than the
preceding one. Let $K$ be a real quadratic field. With the usual notations,
we have $r_{1}=2$ and $r_{2}=0$, so $r=r_{1}+r_{2}-1=1$. The unit theorem
(\hyperref[sec:4.4]{§4}, Theorem \ref{thm:IV.4.1}) implies that the group of units of $K$ is
isomorphic to the product of $\Z$ with the group of roots of unity contained
in $K$. As $K$ admits an imbedding into $\R$, the only roots of unity are
$\pm1$. Thus, assuming that $K$ has been imbedded into $\R$, we have:

\begin{proposition}
\label{prop:IV.6.1}
The positive units of a real quadratic field $K\subset\R$ form a
(multiplicative) group isomorphic to $\Z$.
\end{proposition}

This group contains one and only one generator larger than one, we call it
the \emph{fundamental unit} of $K$.

Let $K=\Q[\sqrt{d}]$, where $d\geq2$ is a square-free integer, and let
$x=a+b\sqrt{d}$ ($a,b\in\Q$) be a unit of $K$. The numbers $x$, $x^{-1}$,
$-x$, and $-x^{-1}$ are units of $K$ and, since
$\mathrm{N}(x)=(a+b\sqrt{d})(a-b\sqrt{d})=\pm1$ (\hyperref[sec:4.4]{§4}, Proposition
\ref{prop:IV.4.1}), these four numbers are $\pm a\pm b\sqrt{d}$. For
$x\neq\pm1$ only one of the four numbers $x,x^{-1},-x,-x^{-1}$ is greater
than one, and it is the largest of the four. Thus the units greater than one
of $K$ are those of the form $a+b\sqrt{d}$ with $a>0$, $b>0$.

(a) Suppose first that $d\equiv2$ or $3\pmod{4}$. In this case the ring of
integers of $K$ is $\Z+\Z\sqrt{d}$ (\hyperref[sec:2.5]{Chapter II, §5}, Theorem
\ref{thm:II.5.1}). As the units of $K$ are integers of norm $\pm1$
(\hyperref[sec:4.4]{§4}, Proposition \ref{prop:IV.4.1}), the units greater than one of $K$ are
the numbers $a+b\sqrt{d}$ with $a,b\in\Z$ and $a>0$, $b>0$ such that
\begin{equation}
a^{2}-db^{2}=\pm1.
\label{eq:IV.6.1}
\end{equation}
We remark that the solutions ``in natural numbers'' $(a,b)$ of equation
\eqref{eq:IV.6.1} (called the ``equation of Pell-Fermat'') are obtained as
follows: take the fundamental unit $a_{1}+b_{1}\sqrt{d}$ of $K$ and put
\begin{equation}
a_{n}+b_{n}\sqrt{d}=(a_{1}+b_{1}\sqrt{d})^{n},\quad n\geq1.
\label{eq:IV.6.2}
\end{equation}
The sequence $(a_{n},b_{n})$ lists all the solutions of \eqref{eq:IV.6.1}.

\begin{statement}{Remarks}
\upshape
\begin{enumerate}[label=(\arabic*),leftmargin=2.0em,itemsep=0.1ex,topsep=0.3ex]
\item It follows from \eqref{eq:IV.6.2} that
$b_{n+1}=a_{1}b_{n}+b_{1}a_{n}$. Since $a_{1},b_{1},a_{n}$ and $b_{n}$ are
all positive, the sequence $(b_{n})$ is strictly increasing. Thus, in order
to explicitly calculate the fundamental unit $a_{1}+b_{1}\sqrt{d}$, it
suffices to write down the sequence $(db^{2})$ for $b\in\N$, $b\geq1$ and to
stop at the first number $db_{1}^{2}$ of this sequence which differs by a
square $a_{1}^{2}$ from $\pm1$. Then $a_{1}+b_{1}\sqrt{d}$ is the
fundamental unit of $K$. For instance, if $d=7$, the sequence $(db^{2})$ is
$7,28,63=64-1=8^{2}-1$, so, taking $b_{1}=3$ and $a_{1}=8$, we see that
$8+3\sqrt{7}$ is the fundamental unit of $\Q[\sqrt{7}]$. We see similarly
that the fundamental units of $\Q[\sqrt{2}]$, $\Q[\sqrt{3}]$, and
$\Q[\sqrt{6}]$ are $1+\sqrt{2}$, $2+\sqrt{3}$, and $5+2\sqrt{6}$. Using the
theory of continued fractions, one can find other, more rapid, procedures
for calculating the fundamental unit.
\item If the fundamental unit is of norm one, the sequence $(a_{n},b_{n})$
gives solutions only for the equation (1$'$) $a^{2}-db^{2}=1$; in this case
the equation (1$''$) $a^{2}-db^{2}=-1$ has no solution in natural numbers.
If the fundamental unit has norm $-1$, the solutions of (1$'$) comprise the
sequence $(a_{2n},b_{2n})$ and those of (1$''$) the sequence
$(a_{2n+1},b_{2n+1})$. The first case occurs when $d=3,6$, or $7$, the
second when $d=2$ or $10$ ($3+\sqrt{10}$ is the fundamental unit in
$\Q[\sqrt{10}]$).
\end{enumerate}
\end{statement}

(b) Assume now that $d\equiv1\pmod{4}$. The integers of $K=\Q[\sqrt{d}]$
are the numbers $\frac{1}{2}(a+b\sqrt{d})$ with $a,b\in\Z$ of the same
parity (\hyperref[sec:2.5]{Chapter II, §5}, Theorem \ref{thm:II.5.1}). Consequently, if
$\frac{1}{2}(a+b\sqrt{d})$ is a unit of $K$ (\hyperref[sec:4.4]{§4}, \hyperref[prop:IV.4.1]{Proposition 2}), we must
have
\begin{equation}
a^{2}-db^{2}=\pm4.
\label{eq:IV.6.3}
\end{equation}
Conversely, if $(a,b)$ is an integer solution of \eqref{eq:IV.6.3}, then
$\frac{1}{2}(a+b\sqrt{d})$ is an integer of $K$ (its trace is $a$ and its
norm, by \eqref{eq:IV.6.3}, is $\pm1$) and, hence, a unit of $K$. As in
(a), writing $a_{1}+b_{1}\sqrt{d}$ for the fundamental unit of $K$, we see
that the solutions in pairs of natural numbers $(a,b)$ of \eqref{eq:IV.6.3}
comprise the values of the sequence $(a_{n},b_{n})$ ($n\geq1$) defined by
setting
\begin{equation}
a_{n}+b_{n}\sqrt{d}=2^{1-n}(a_{1}+b_{1}\sqrt{d})^{n}.
\label{eq:IV.6.4}
\end{equation}
The calculation of $a_{1}+b_{1}\sqrt{d}$ may be accomplished as in (a). For
example, the fundamental units of $\Q[\sqrt{5}]$, $\Q[\sqrt{13}]$, and
$\Q[\sqrt{17}]$ are $\frac{1}{2}(1+\sqrt{5})$, $\frac{1}{2}(3+\sqrt{13})$,
and $4+\sqrt{17}$; these three units all have norm $-1$. For the choice of
the sign $\pm1$ in \eqref{eq:IV.6.3} we have results similar to those
obtained in (a).

\begin{statement}{Remark}
\upshape
In the case $d\equiv1\pmod{4}$ the solutions of the Pell-Fermat equation
\begin{equation}
a^{2}-db^{2}=\pm1
\label{eq:IV.6.5}
\end{equation}
correspond to units $a+b\sqrt{d}$ ($a,b>0$) which belong to the ring
$B=\Z[\sqrt{d}]$. This ring $B$ is a subring of the ring $A$ of integers of
$K$ and the positive units of $B$ form a subgroup $G$ of the group of
positive units of $A$. Let $u=\frac{1}{2}(a+b\sqrt{d})$ be the fundamental
unit of $K$. If $a$ and $b$ are both even, then $u\in B$, so that $G$
consists of the powers of $u$ (this is the case, for instance, when $d=17$).
If $a$ and $b$ are both odd, then $u^{3}\in B$. (To see this note that
$8u^{3}=a(a^{2}+3b^{2}d)+b(3a^{2}+b^{2}d)\sqrt{d}$. Since
$a^{2}-b^{2}d=\pm4$, $a^{2}+3b^{2}d=4(b^{2}d\pm1)$, which is a multiple of
$8$, since $b$ and $d$ are odd. Similarly $3a^{2}+b^{2}d=4(a^{2}\pm1)$,
which is again, because $a$ is odd, a multiple of $8$.) In this case $G$ is
comprised of the powers of $u^{3}$ ($u^{2}\notin B$, otherwise
$u=u^{3}/u^{2}\in B$). This happens, for instance, when $d=5$
(respectively, $d=13$), in which case $u^{3}=2+\sqrt{5}$ (respectively,
$u^{3}=18+5\sqrt{13}$).
\end{statement}

%% ----------------------------------------------------------------------
\sechead{7}{A generalization of the unit theorem}%
        {A GENERALIZATION OF THE UNIT THEOREM}

\begin{proposition}
\label{prop:IV.7.1}
Let $A$ be a ring which is a $\Z$-module of finite type. Then the
multiplicative group $A^{*}$ consisting of the units of $A$ is a
(commutative) group of finite type.
\end{proposition}

\begin{statement}{Remarks}
\upshape
For a commutative group $G$, ``of finite type'' means ``of finite type with
respect to the $\Z$-module structure of $G$''. A subgroup of a commutative
group of finite type is of finite type (\hyperref[sec:3.1]{Chapter III, §1}, Corollary
\ref{cor:III.1.2} of Theorem \ref{thm:III.1.1}). Let us note that $A$ is a
Noetherian ring, for the ideals of $A$ are $\Z$-submodules of $A$.
\end{statement}

\begin{proof}[Proof of the proposition]
First we treat the case when $A$ is an integral domain. If its field of
fractions $K$ is of characteristic zero, $K$ is a finite-dimensional vector
space over $\Q$, so $K$ is a number field. On the other hand, $A$ is integral
over $\Z$ (since it is a $\Z$-module of finite type, cf.\ \hyperref[sec:2.1]{Chapter II, §1},
Theorem \ref{thm:II.1.1}), and, therefore, $A$ is a subring of $B$, the ring
of integers of $K$. Thus $A\subset B$ and, since $B$ is of finite type by the
unit theorem (\hyperref[sec:4.4]{§4}, Theorem \ref{thm:IV.4.1}), so is $A$. If $K$ is of
characteristic $p\neq0$, $K$ is a finite extension of $\F_{p}$, so $K$ is a
finite field and $A^{*}$ is a finite group.

Now let us consider the case in which $A$ is \emph{reduced} (by definition,
this means that $0$ is the only nilpotent element in $A$). We shall need the
following lemma.
\end{proof}

\begin{statement}{Lemma}
In a reduced Noetherian ring $A$, the ideal $(0)$ is expressible as the
intersection of finitely many prime ideals.
\end{statement}

\begin{proof}
We know that, in a Noetherian ring, any ideal contains a product of prime
ideals (\hyperref[sec:3.3]{Chapter III, §3}, \hyperref[lem:III.3.3]{Lemma 3}). $(0)$ is the smallest ideal, so $(0)$ is a
product of prime ideals
$(0)=\mathfrak{p}_{1}^{n_{1}}\cdots\mathfrak{p}_{q}^{n_{q}}$. Let
$x\in\mathfrak{p}_{1}\cap\cdots\cap\mathfrak{p}_{q}$. Then
$x^{n_{1}+\cdots+n_{q}}\in
\mathfrak{p}_{1}^{n_{1}}\cdots\mathfrak{p}_{q}^{n_{q}}=(0)$, so
$x^{n_{1}+\cdots+n_{q}}=0$. Since $A$ is reduced, this means that $x=0$.
Therefore, $\mathfrak{p}_{1}\cap\cdots\cap\mathfrak{p}_{q}=(0)$.
\end{proof}

Now, returning to the proof of Proposition \ref{prop:IV.7.1}, we let
$(0)=\mathfrak{p}_{1}\cap\cdots\cap\mathfrak{p}_{q}$, the $\mathfrak{p}_{i}$
being prime ideals. It follows that the canonical homomorphism
$\varphi:A\to\prod_{i=1}^{q}A/\mathfrak{p}_{i}$ is injective. An element of a
ring product is invertible if and only if all of its components are
invertible, so that
$(\prod_{i}A/\mathfrak{p}_{i})^{*}=\prod_{i}(A/\mathfrak{p}_{i})^{*}$. By
the integral domain case, each $(A/\mathfrak{p}_{i})^{*}$ is of finite type;
therefore, $\prod(A/\mathfrak{p}_{i})^{*}$ is of finite type and so is any
subgroup, e.g.\ $\varphi(A^{*})$ (recall that $\Z$ is a Noetherian ring).
Since $\varphi$ is injective, $A^{*}$ is of finite type.

Let us finally consider the general case. Observe that the set
$\mathfrak{n}$ of nilpotent elements of $A$ is an ideal, since $x^{p}=0$,
$y^{q}=0$, and $a\in A$ imply that $(x+y)^{p+q-1}=0$ and $(ax)^{p}=0$. On
the other hand there exists an integer $s$ such that
$\mathfrak{n}^{s}=(0)$. (To see this, let $(x_{1},\ldots,x_{r})$ be a finite
set of generators of the ideal $\mathfrak{n}$ in the Noetherian ring $A$.
Assume $x_{i}^{q_{i}}=0$, $i=1,\ldots,r$. Then, with
$s=q_{1}+\cdots+q_{r}$, it is clear that $\mathfrak{n}^{s}=0$.) We proceed by
induction on $s$. The case $s=1$ is the reduced case---already treated.
Suppose $s>1$ and write $\varphi$ for the canonical homomorphism
$\varphi:A\to A/\mathfrak{n}^{s-1}$. Then
$\varphi(A^{*})\subset(A/\mathfrak{n}^{s-1})^{*}$, so $\varphi(A^{*})$ is of
finite type. The kernel of $\varphi$ restricted to $A^{*}$ belongs to
$1+\mathfrak{n}^{s-1}$; in fact,
$\ker(\varphi)=1+\mathfrak{n}^{s-1}$, because for $s>1$,
$(\mathfrak{n}^{s-1})^{2}\subset\mathfrak{n}^{s}=(0)$, which implies that
any element $1+x\in1+\mathfrak{n}^{s-1}$ has $1-x$ as its inverse:
$(1+x)(1-x)=1-x^{2}=1$. To prove that $1+\mathfrak{n}^{s-1}$ is of finite
type we need only observe that the mapping $x\mapsto1+x$ of the additive
group $\mathfrak{n}^{s-1}$ to $1+\mathfrak{n}^{s-1}$ is an isomorphism
($\mathfrak{n}^{s-1}$ is of finite type because it is a submodule of $A$).
But this is obvious from the relation $(1+x)(1+y)=1+x+y$ for
$x,y\in\mathfrak{n}^{s-1}$. That $A^{*}$ is of finite type now follows from
\hyperref[sec:3.1]{Chapter III, §1}, Proposition \ref{prop:III.1.1}.

Using arguments borrowed from algebraic geometry, one can show that, for any
\emph{reduced} ring $B$ of the form $B=\Z[x_{1},\ldots,x_{n}]$
(i.e.\ finitely generated over $\Z$ as a ring), the group $B^{*}$ of units
in $B$ is of finite type (\hyperref[bib:6]{[6]}).

%% ----------------------------------------------------------------------
\clearpage
\annexhead{The calculation of a volume}{THE CALCULATION OF A VOLUME}

\begin{statement}{Proposition}
Let $r_{1}$ and $r_{2}\in\N$, $n=r_{1}+2r_{2}$, $t\in\R$, and let $B_{t}$ be
the set of all elements
$(y_{1},\ldots,y_{r_{1}},z_{1},\ldots,z_{r_{2}})\in
\R^{r_{1}}\times\C^{r_{2}}$ such that
\begin{equation}
\sum_{i=1}^{r_{1}}|y_{i}|+2\sum_{j=1}^{r_{2}}|z_{j}|\leq t.
\label{eq:IV.A.1}
\end{equation}
Let $\mu$ denote the Lebesgue measure in $\R^{n}$. Then,
\begin{equation}
\mu(B_{t})=2^{r_{1}}\Bigl(\frac{\pi}{2}\Bigr)^{r_{2}}\frac{t^{n}}{n!}
\quad\text{for any }t\geq0.
\label{eq:IV.A.2}
\end{equation}
\end{statement}

\begin{proof}
We set $\mu(B_{t})=V(r_{1},r_{2},t)$ and argue by double induction on $r_{1}$
and $r_{2}$. Clearly $V(1,0,t)=2t$ (the segment $[-t,+t]$) and
$V(0,1,t)=\pi t^{2}/4$ (the disc of radius $t/2$). These results verify
\eqref{eq:IV.A.2} in the special cases considered.

Now assume \eqref{eq:IV.A.2} gives $V(r_{1},r_{2},t)$. First we compute
$V(r_{1}+1,r_{2},t)$. The set
$B_{t}\subset\R\times\R^{r_{1}}\times\C^{r_{2}}$, which corresponds to
$r_{1}+1$ and $r_{2}$, is defined by the relation
\[
|y|+\sum_{i=1}^{r_{1}}|y_{i}|+2\sum_{j=1}^{r_{2}}|z_{j}|\leq t
\quad\text{with }y\in\R.
\]
Integrating ``in strips'' and observing that for $|y|>t$,
$B_{t}=\emptyset$, we see that
\[
V(r_{1}+1,r_{2},t)=\int_{-t}^{+t}V(r_{1},r_{2},t-|y|)\,dy.
\]
Use of the induction hypothesis gives
\[
V(r_{1}+1,r_{2},t)=2\int_{0}^{t}2^{r_{1}}
\Bigl(\frac{\pi}{2}\Bigr)^{r_{2}}\frac{(t-y)^{n}}{n!}\,dy
=2^{r_{1}+1}\Bigl(\frac{\pi}{2}\Bigr)^{r_{2}}
\frac{t^{n+1}}{(n+1)!},
\]
which agrees with \eqref{eq:IV.A.2}.

It remains to compute $V(r_{1},r_{2}+1,t)$. We keep the induction hypothesis
that $V(r_{1},r_{2},t)$ satisfies \eqref{eq:IV.A.2}. The set
$B_{t}\subset\R^{r_{1}}\times\C^{r_{2}}\times\C$ which corresponds to
$r_{1}$ and $r_{2}+1$ is defined by the relation
\[
\sum_{i=1}^{r_{1}}|y_{i}|+2\sum_{j=1}^{r_{2}}|z_{j}|+2|z|\leq t
\quad\text{with }z\in\C.
\]
Again integrating in strips, we obtain
\[
V(r_{1},r_{2}+1,t)=\int_{|z|\leq t/2}
V(r_{1},r_{2},t-2|z|)\,d\mu(z),
\]
where $d\mu(z)$ denotes Lebesgue measure on $\C$. Putting
$z=\rho e^{i\theta}$ ($\rho\in\R_{+}$, $0\leq\theta\leq2\pi$), we have
$d\mu(z)=\rho\,d\rho\,d\theta$. Use of the induction hypothesis gives
\[
\begin{aligned}
V(r_{1},r_{2}+1,t)
&=\int_{0}^{t/2}\int_{0}^{2\pi}2^{r_{1}}
\Bigl(\frac{\pi}{2}\Bigr)^{r_{2}}\frac{(t-2\rho)^{n}}{n!}
\rho\,d\rho\,d\theta\\
&=2^{r_{1}}\Bigl(\frac{\pi}{2}\Bigr)^{r_{2}}\frac{2\pi}{n!}
\int_{0}^{t/2}(t-2\rho)^{n}\rho\,d\rho.
\end{aligned}
\]
Calculating $\int_{0}^{t/2}(t-2\rho)^{n}\rho\,d\rho$ by putting $2\rho=x$
and integrating by parts, we find that this integral has the value
$t^{n+2}/[4(n+1)(n+2)]$. Thus,
\[
V(r_{1},r_{2}+1,t)=2^{r_{1}}\Bigl(\frac{\pi}{2}\Bigr)^{r_{2}+1}
\frac{t^{n+2}}{(n+2)!},
\]
which agrees with \eqref{eq:IV.A.2} since $r_{1}+2(r_{2}+1)=n+2$.
\end{proof}
